\begin{trapbox}

	\begin{description}[style=nextline, leftmargin=2.8em, labelsep=0.5em, itemsep=0.35em]
		\item[\textbf{\textcolor{CTrap}{易错点一（方向混淆）}}]
			常见错误：看到$A\subseteq B$就认为$p$是$q$的必要条件。
		\item[\textbf{\textcolor{CTrap}{正确理解（集小推充分）：}}]
			若$A\subseteq B$，则$p\Rightarrow q$，即$p$是$q$的充分条件。
		\item[\textbf{\textcolor{CTrap}{错因分析（记反定义）：}}]
			混淆了充分条件与必要条件的定义方向。
	\end{description}

	\medskip
	\begin{description}[style=nextline, leftmargin=2.8em, labelsep=0.5em, itemsep=0.35em]
		\item[\textbf{\textcolor{CTrap}{易错点二（端点比较遗漏）}}]
			比较数集包含关系时只检查了一个端点，例如由$[-1,3]\subseteq[-m,m]$只得到$m\ge 3$而遗漏$-m\le -1$。
		\item[\textbf{\textcolor{CTrap}{正确理解（两端都满足）：}}]
			需同时满足左端点和右端点的不等式：$-m\le -1$且$m\ge 3$。
		\item[\textbf{\textcolor{CTrap}{错因分析（只考虑一端）：}}]
			思维不全面，忽视区间左右边界。
	\end{description}

	\medskip
	\begin{description}[style=nextline, leftmargin=2.8em, labelsep=0.5em, itemsep=0.35em]
		\item[\textbf{\textcolor{CTrap}{易错点三（忽略相等判断）}}]
			当$A=B$时，错误地认为是充分不必要条件。
		\item[\textbf{\textcolor{CTrap}{正确理解（相等即充要）：}}]
			若$A=B$，则$p\Leftrightarrow q$，为充要条件。
		\item[\textbf{\textcolor{CTrap}{错因分析（未核查相等）：}}]
			只验证了$A\subseteq B$，没有检查$B\subseteq A$是否也成立。
	\end{description}

\end{trapbox}
